\Chapter{Mathematical Logic}
\Lettrine{Mathematics} is an inherently logical subject; the basis of all mathematics is in how we form logical statements and prove them. For one to begin with the exploration of axioms in mathematics, we have to start with the fundamentals -- mathematical logic.

\section{Mathematical Statements}
The core of mathematics is in writing mathematical statements. Maybe your mathematics experience so far has mostly involved finding answers to problems. However, as we begin defining numbers from scratch, we must be adept at writing and communicating mathematical ideas. Precision is vital, so we must first consider the basic building blocks of all mathematical ideas: mathematical statements.

\begin{definition}[Statements]
    A \term{statement}\index{statement} (or \term{proposition}\index{proposition}) is a sentence that is definitely true or definitely false, but not both.
\end{definition}

In other words, a mathematical statement is a sentence we can prove or disprove. A mathematical statement is not an opinion, not an order, not an instruction, nor a feeling. A statement is either a fact or falsity. There is no alternate view that one can take to prove otherwise.

\begin{example}
    The following are statements.
    \begin{itemize}
        \item A square is a quadrilateral.
        \item 16 is a square number.
        \item 9 is a prime.
        \item Every even number greater than 2 can be expressed as the sum of two primes.
        \item $3 + 8 = 12$.
    \end{itemize}

    \newpage

    The following are \textbf{not} statements.
    \begin{itemize}
        \item Would you like some water?
        \item The product of two numbers.
        \item Shoo!
        \item $1 + 2 + 3 + \cdots + n$.
        \item $x + 1 = 3$.
    \end{itemize}
\end{example}

The reason why ``$x + 1 = 3$'' is not a statement because its truth or falsity depends on the value of $x$. For example, if $x = 2$ then the statement is true, but if $x$ is any other value the statement is false. We will examine these kinds of sentences later.

We may name statements using lowercase variables like $p$, $q$, and $r$.

\begin{example}
    If we define the statements $p$, $q$, and $r$, where
    \begin{align*}
        p: &\ \text{every odd number is one more than an even number},\\
        q: &\ \text{every triangle has sides of equal length}, \text{ and}\\
        r: &\ \text{the only even prime number is 2},
    \end{align*}
    then $p$ is a true statement, $q$ is a false statement, and $r$ is a true statement.
\end{example}

\begin{exercise}[Statement Detective]
    Which of the following are statements?

    \begin{partquestions}{\alph*}[2]
        \item Hello!
        \item 4 equals 5.
        \item You passed the exam.
        \item $n \geq 4$.
        \item $8 = 2 \times 4$.
        \item There is a number $n$ that is one less than a square.
        \item $x^2 - 1 = (x-1)(x+1)$.
    \end{partquestions}
\end{exercise}

\section{Logical Connectives}
Statements come in many different forms. We took a look at some examples and non-examples of statements above. Some statements can be broken down into `simpler' statements, while others cannot. We have names for these kinds of statements.

\begin{definition}[Atomic and Molecular Statements]
    An \term{atomic statement}\index{statement!atomic} is a statement that cannot be divided into other simpler statements. A statement that is not atomic is called a \term{molecular statement}\index{statement!molecular}.
\end{definition}

\begin{example}
    \vspace{-4mm}
    \begin{multicols}{2}
        {
            These are atomic statements.
            \begin{itemize}
                \item A square is a quadrilateral.
                \item 16 is a square number.
                \item 9 is prime.
                \item $3 + 8 = 12$.
            \end{itemize}
        }
        {
            These are molecular statements.
            \begin{itemize}
                \item $1 + 1 = 2$ and $2 + 2 = 4$.
                \item Either 8 is prime or not.
                \item 18 is a multiple of 3 and 6.
                \item 10 is even or odd.
            \end{itemize}
        }
    \end{multicols}
\end{example}

To form new more complex (i.e., molecular) statements from simpler (i.e., atomic or molecular) statements, we can use \term{logical connectives}\index{logical connective}. These connectives, as implied by their name, connect statements together to form a more complex statement. To define these logical connectives, we elect to use truth tables.

\begin{definition}
    A \term{truth table}\index{truth table} is a table showing how the truth or falsity of a statement depends on the truth or falsity of the statements from which it is constructed from.
\end{definition}

This definition will seem a little out of place without examples, so let's define a few logical connectives.

\subsection{Negation}
Negation is probably the easiest logical connective to discuss. Negation simply turns any true statement into a false statement, and vice versa.

\begin{example}
    The negation of ``5 is an odd number'', which is true, is ``5 is \textbf{not} an odd number'', i.e. ``5 is an even number'', which is false.
\end{example}

We can capture the idea of negation in the following definition.

\begin{definition}[Negation]
    Let $p$ be a statement. The \term{negation of $p$}\index{logical connective!negation}, denoted $\lnot p$ and read ``not $p$'', has the following truth table.
    \begin{table}[H]
        \centering
        \begin{tabular}{|c||c|}
            \hline
            $\boldsymbol{p}$ & $\boldsymbol{\lnot p}$ \\ \hhline{|=#=|}
            T & F \\ \hline
            F & T \\ \hline
        \end{tabular}
    \end{table}
\end{definition}

We read the truth table as follows.
\begin{itemize}
    \item The leftmost column(s) show the `source' statements. In the above definition, we only have $p$ as the singular `source' statement.
    \item The rightmost column show the `output' statement. In the above definition, we have $\lnot p$ as the `output' statement.
    \item We denote the possibility that a statement is true by ``T'' and denote the possibility that a statement is false by ``F''.
\end{itemize}

Thus, the above table tells us that $\lnot p$ is false when $p$ is true and $\lnot p$ is true when $p$ is false.

\begin{example}
    Let $p$ be the statement ``5 is an odd number'', which is a true statement. Then $\lnot p$ is ``5 is \textbf{not} an odd number'' (i.e., ``5 is an even number''), which is a false statement.
\end{example}

\subsection{Conjunction}
The conjunction connective is better remembered as ``AND''.

\begin{definition}[Conjunction]
    Let $p$ and $q$ be statements. The \term{conjunction of $p$ and $q$}\index{logical connective!conjunction}, denoted $p \land q$ and read ``$p$ and $q$'', has the following truth table.
    \begin{table}[H]
        \centering
        \begin{tabular}{|c|c||c|}
            \hline
            $\boldsymbol{p}$ & $\boldsymbol{q}$ & $\boldsymbol{p \land q}$ \\ \hhline{|=|=#=|}
            T & T & T \\ \hline
            T & F & F \\ \hline
            F & T & F \\ \hline
            F & F & F \\ \hline
        \end{tabular}
    \end{table}
\end{definition}

In other words, the conjunction of two statements is true if and only if both the statements are true.

\begin{example}
    We look at a series of statements and their conjunction.
    \begin{itemize}
        \item If $p$ is the statement ``3 is odd'' (true) and $q$ is the statement ``4 is even'' (true), then $p \land q$ is the statement ``3 is odd \textbf{and} 4 is even'' (true).
        \item If $p$ is the statement ``a square is a quadrilateral'' (true) and $q$ is the statement ``a triangle is a quadrilateral'' (false), then $p \land q$ is the statement ``a square is a quadrilateral \textbf{and} a triangle is a quadrilateral'' (false).
        \begin{itemize}
            \item Alternatively, $p \land q$ can be written as ``both squares and triangles are quadrilaterals'' which is still false.
        \end{itemize}
        \item If $p$ is the statement ``4 equals 5'' (false) and $q$ is the statement ``5 is odd'' (true), then $p \land q$ is the statement ``4 equals 5 \textbf{and} 5 is odd'' (false).
        \item If $p$ is the statement ``9 is a prime'' (false) and $q$ is the statement ``$3 + 8 = 12$'' (false), then $p \land q$ is the statement ``9 is prime \textbf{and} $3 + 8 = 12$'' (false).
    \end{itemize}
\end{example}

\subsection{Disjunction}
Like with the conjunction connective, the word ``disjunction'' is not particularly enlightening. The disjunction connective is thus better understood to be akin to ``OR''.

\begin{definition}[Disjunction]
    Let $p$ and $q$ be statements. The \term{disjunction of $p$ and $q$}\index{logical connective!disjunction}, denoted $p \lor q$ and read ``$p$ or $q$'', has the following truth table.
    \begin{table}[H]
        \centering
        \begin{tabular}{|c|c||c|}
            \hline
            $\boldsymbol{p}$ & $\boldsymbol{q}$ & $\boldsymbol{p \lor q}$ \\ \hhline{|=|=#=|}
            T & T & T \\ \hline
            T & F & T \\ \hline
            F & T & T \\ \hline
            F & F & F \\ \hline
        \end{tabular}
    \end{table}
\end{definition}

Thus, the disjunction of two statements is true if and only if at least one of the statements is true.

\begin{example}
    We look at a series of statements and their disjunction.
    \begin{itemize}
        \item If $p$ is the statement ``3 is odd'' (true) and $q$ is the statement ``4 is even'' (true), then $p \lor q$ is the statement ``3 is odd \textbf{or} 4 is even'' (true).
        \item If $p$ is the statement ``a square is a quadrilateral'' (true) and $q$ is the statement ``a triangle is a quadrilateral'' (false), then $p \lor q$ is the statement ``a square is a quadrilateral \textbf{or} a triangle is a quadrilateral'' (true).
        \item If $p$ is the statement ``4 equals 5'' (false) and $q$ is the statement ``5 is odd'' (true), then $p \lor q$ is the statement ``4 equals 5 \textbf{or} 5 is odd'' (true).
        \item If $p$ is the statement ``9 is a prime'' (false) and $q$ is the statement ``$3 + 8 = 12$'' (false), then $p \lor q$ is the statement ``9 is prime \textbf{or} $3 + 8 = 12$'' (false).
    \end{itemize}
\end{example}

\begin{exercise}[Combining Basic Connectives]
    Define the statements $p$, $q$, and $r$, where
    \begin{align*}
        p: &\ \text{1 is a positive number},\\
        q: &\ \text{1 is an even number}, \text{ and}\\
        r: &\ \text{1 is an odd number}.
    \end{align*}
    Is the statement ``$\lnot((p \lor q) \land r)$ is false'' true?
\end{exercise}

\subsection{Conditional/Implication}
The conditional (or implication) connective is harder to motivate. Connecting the statements $p$ and $q$ with the conditional has essentially the same meaning as ``if $p$ then $q$''. However, there is a notable quirk with the definition of the conditional connective which we will see below.

\begin{definition}[Conditional/Implication]
    Let $p$ and $q$ be statements. The statement connecting $p$ and $q$ with the \term{conditional}\index{logical connective!conditional} or \term{implication}\index{logical connective!implication} connective, denoted $p \implies q$ and typically read ``$p$ implies $q$'', has the following truth table.
    \begin{table}[H]
        \centering
        \begin{tabular}{|c|c||c|}
            \hline
            $\boldsymbol{p}$ & $\boldsymbol{q}$ & $\boldsymbol{p \implies q}$ \\ \hhline{|=|=#=|}
            T & T & T \\ \hline
            T & F & F \\ \hline
            F & T & T \\ \hline
            F & F & T \\ \hline
        \end{tabular}
    \end{table}
    A statement of the form ``$p$ implies $q$'' is called a \term{conditional statement}\index{statement!conditional}, and $p$ is called the \term{antecedent}\index{antecedent} while $q$ is called the \term{consequent}\index{consequent}.
\end{definition}

\begin{remark}
    The statement ``$p \implies q$'' can be read many ways. For instance,
    \begin{itemize}
        \item $p$ implies $q$;
        \item if $p$ then $q$;
        \item $q$ if $p$;
        \item $p$ only if $q$;
        \item $p$ is a sufficient condition for $q$; and
        \item $q$ is a necessary condition for $p$.
    \end{itemize}
\end{remark}

Let us motivate the reason why we define the truth table for the conditional as above. Suppose a mother tells her child that ``if it is sunny, then we will go to the park''. In this case, the statement $p$ is ``it is sunny'' and the statement $q$ is ``we will go to the park'', making the overall statement $p \implies q$. So what are the possibilities for the truth or falsity of this overall statement?
\begin{itemize}
    \item If it is sunny and they went to the park, the promise of the overall statement was not broken. Thus, if $p$ and $q$ are both true then $p \implies q$ is also true.
    \item If it is sunny and they did not go to the park, clearly the mother broke her promise. So if $p$ is true but $q$ is false then we can say that $p \implies q$ is false.
    \item If instead it is not sunny (maybe it is cloudy) and they went to the park, at no point was the promise broken. The only thing that the mother said was that they would go to the park if it was sunny; it was not the \textbf{only} reason that they could go to the park. Perhaps the child was very well behaved and she wanted to reward her child. Perhaps they were on the way back from school and coincidentally pass by the park. Either way, if $p$ is false but $q$ is true, we say that $p \implies q$ is true.
    \item If it is not sunny and they did not go to the park, it is intuitive to say that the promise was not broken. It could be raining, and so going to the park would be a terrible idea! So if both $p$ and $q$ are false then $p \implies q$ is still true.
\end{itemize}

Thus, the only case when $p \implies q$ is false is if $p$ is true but $q$ is false.

\begin{remark}
    A conditional statement that is true by virtue of the antecedent being false is said to be \term{vacuously true}\index{vacuously true}.
\end{remark}

\begin{example}
    \begin{itemize}[leftmargin=*]
        \item The statement ``if 3 is odd then 4 is even'' is clearly true, with both antecedent and consequent being true.
        \item The statement ``if 5 is prime then 6 is prime'' is clearly false, with a true antecedent and a false consequent.
        \item The statement ``if 0 equals 1 then 5 is odd'' is true since the antecedent is false, meaning that the statement is vacuously true (with a true consequent).
        \item The statement ``if 1 equals 2 then 1234 equals 5678'' is true since the antecedent is false, meaning that this is a vacuously true statement (with a false consequent).
    \end{itemize}    
\end{example}

\subsection{Biconditional}
Sometimes we would want to form a statement of the form ``if $p$ then $q$, and if $p$ then $q$''. Using symbolic notation we could write it as ``$(p \implies q) \land (q \implies p)$'', but this is a little clunky to write every time we wish to write such a statement. This motivates the biconditional connective.

\begin{definition}[Biconditional]
    Let $p$ and $q$ be statements. The statement connecting $p$ and $q$ with the \term{biconditional}\index{logical connective!biconditional} connective, denoted $p \iff q$ and typically read ``$p$ if and only if $q$'', has the following truth table.
    \begin{table}[H]
        \centering
        \begin{tabular}{|c|c||c|}
            \hline
            $\boldsymbol{p}$ & $\boldsymbol{q}$ & $\boldsymbol{p \iff q}$ \\ \hhline{|=|=#=|}
            T & T & T \\ \hline
            T & F & F \\ \hline
            F & T & F \\ \hline
            F & F & T \\ \hline
        \end{tabular}
    \end{table}
    A statement of the form ``$p$ if and only if $q$'' is called a \term{biconditional statement}\index{statement!biconditional}.
\end{definition}

\begin{remark}
    The statement ``$p \iff q$'' can also be read many ways. For instance,
    \begin{itemize}
        \item $p$ if and only if $q$;
        \item $p$ iff $q$ (i.e., ``iff'' is a shorthand for ``if and only if'');
        \item $p$ is a necessary and sufficient condition for $q$; and
        \item $p$ is equivalent to $q$.
    \end{itemize}
\end{remark}

\begin{exercise}[(Bi)conditionals]
    Let $t$ be the conditional statement ``20 is a multiple of 5 if 20 ends in a 0 or 5''. Let $r$ be the antecedent of $t$ and $s$ be the consequent of $t$.
    \begin{partquestions}{\roman*}
        \item Write down the statement $s$ in English.
        \item Define atomic statements $p$ and $q$ such that $r$ is a molecular statement formed by $p$ and $q$.
        \item Write the statement $t$ using only atomic statements and logical connectives.
        \item Write the biconditional statement ``$s \iff r$'' in English. Is ``$s \iff r$'' true?
    \end{partquestions}
\end{exercise}

\section{Logical Equivalence}
Sometimes, the molecular statements that we are considering will have an equivalent form that is simpler to work with. This means that, regardless of the underlying truth values of the atomic statements making up the molecular statements, they hold the same truth value. This is the idea of logical equivalences.

\begin{definition}
    Two statements $p$ and $q$ are \term{logically equivalent}\index{logically equivalent} if and only if they have identical truth values, regardless of the atomic statements that make up them.

    We write $p \equiv q$ if and only if the statements $p$ and $q$ are logically equivalent.
\end{definition}

\begin{remark}[For the Pedantic]
    Logical equivalences is really the same as biconditionals. That is, when we write ``$p \equiv q$'' it is the same as ``$p \iff q$''. However, for this book, we will write ``$p \equiv q$'' to mean ``$p$ and $q$ have the same truth table'', whereas we will write ``$p \iff q$'' if we intend to mean that ``if $p$ then $q$ and if $q$ then $p$'' (i.e., the `traditional' biconditional meaning).
\end{remark}

To show that two statements are logically equivalent, we typically consider their truth tables.

\begin{example}
    Let $p$ and $q$ be statements. We show that $(p \implies q) \equiv (\lnot p) \lor q$ by considering the truth tables of the two statements.
    \begin{table}[H]
        \centering
        \begin{tabular}{|c|c||c||c|c|}
            \hline
            $\boldsymbol{p}$ & $\boldsymbol{q}$ & $\boldsymbol{\lnot p}$ & $\boldsymbol{p \implies q}$ & $\boldsymbol{(\lnot p) \lor q}$ \\ \hhline{|=|=#=#=|=|}
            T & T & F & T & T \\ \hline
            T & F & F & F & F \\ \hline
            F & T & T & T & T \\ \hline
            F & F & T & T & T \\ \hline
        \end{tabular}
    \end{table}
    By inspection, we can see that $(\lnot p) \lor q$ has the same truth values as $p \implies q$ regardless of the truth values of $p$ and $q$. Thus, the $p \implies q$ is logically equivalent to $(\lnot p) \lor q$.
\end{example}
\begin{remark}
    We separate the intermediate value(s) from the rest by drawing a double line to the sides of the intermediate values.
\end{remark}

\begin{example}
    We show that
    \[
        (p \iff q) \equiv (p \land q) \lor ((\lnot p) \land (\lnot q))
    \]
    for any statements $p$ and $q$. For brevity, let $r = p \land q$, $s = (\lnot p) \land (\lnot q)$, and $t = (p \land q) \lor ((\lnot p) \land (\lnot q))$. So we are to show that $(p \iff q) \equiv t$.
    \begin{table}[H]
        \centering
        \begin{tabular}{|c|c||c|c|c|c||c|c|}
            \hline
            $\boldsymbol{p}$ & $\boldsymbol{q}$ & $\boldsymbol{\lnot p}$ & $\boldsymbol{\lnot q}$ & $\boldsymbol{r}$ & $\boldsymbol{s}$ & $\boldsymbol{p \iff q}$ & $\boldsymbol{t}$ \\ \hhline{|=|=#=|=|=|=#=|=|}
            T & T & F & F & T & F & T & T \\ \hline
            T & F & F & T & F & F & F & F \\ \hline
            F & T & T & F & F & F & F & F \\ \hline
            F & F & T & T & F & T & T & T \\ \hline
        \end{tabular}
    \end{table}
    By inspection of the truth table we establish the required result.
\end{example}

\begin{exercise}\label{exercise:contrapositive}
    Show that
    \[
        (p \implies q) \equiv ((\lnot q) \implies (\lnot p))
    \]
    for all statements $p$ and $q$. 
\end{exercise}

What \myref{exercise:contrapositive} reveals is an important relationship regarding conditional statements.

\begin{definition}[Contrapositive]
    Let $p$ and $q$ be statements. The \term{contrapositive}\index{contrapositive} of the conditional statement $p \implies q$ is $(\lnot q) \implies (\lnot p)$.
\end{definition}
\begin{proposition}
    The contrapositive of a conditional statement is equivalent to the conditional statement.
\end{proposition}
\begin{proof}
    See \myref{exercise:contrapositive}.
\end{proof}

There are other properties of logical connectives that we will note down in the following propositions.

\begin{proposition}[Commutativity of AND and OR]
    Let $p$ and $q$ be statements. Then $p \land q \equiv q \land p$ and $p \lor q \equiv q \lor p$.
\end{proposition}
\begin{proof}
    We consider truth tables.
    \begin{table}[H]
        \centering
        \begin{minipage}{.45\linewidth}
            \centering
            \begin{tabular}{|c|c||c|c|}
                \hline
                $\boldsymbol{p}$ & $\boldsymbol{q}$ & $\boldsymbol{p \land q}$ & $\boldsymbol{q \land p}$ \\ \hhline{|=|=#=|=|}
                T & T & T & T \\ \hline
                T & F & F & F \\ \hline
                F & T & F & F \\ \hline
                F & F & F & F \\ \hline
            \end{tabular}
        \end{minipage}
        \begin{minipage}{.45\linewidth}
            \centering
            \begin{tabular}{|c|c||c|c|}
                \hline
                $\boldsymbol{p}$ & $\boldsymbol{q}$ & $\boldsymbol{p \lor q}$ & $\boldsymbol{q \lor p}$ \\ \hhline{|=|=#=|=|}
                T & T & T & T \\ \hline
                T & F & T & T \\ \hline
                F & T & T & T \\ \hline
                F & F & F & F \\ \hline
            \end{tabular}
        \end{minipage}
    \end{table}
    From the above truth tables, the result is established.
\end{proof}

\begin{proposition}[Associativity of AND and OR]
    Let $p$, $q$, and $r$ be statements. Then $p \land (q \land r) \equiv (p \land q) \land r$ and $p \lor (q \lor r) \equiv (p \lor q) \lor r$.
\end{proposition}
\begin{proof}
    We first show $p \land (q \land r) \equiv (p \land q) \land r$.
    \begin{table}[H]
        \centering
        \begin{tabular}{|c|c|c||c|c||c|c|}
            \hline
            $\boldsymbol{p}$ & $\boldsymbol{q}$ & $\boldsymbol{r}$ & $\boldsymbol{p \land q}$ & $\boldsymbol{q \land r}$ & $\boldsymbol{p \land (q \land r)}$ & $\boldsymbol{(p \land q) \land r}$ \\ \hhline{|=|=|=#=|=#=|=|}
            T & T & T & T & T & T & T \\ \hline
            T & T & F & T & F & F & F \\ \hline
            T & F & T & F & F & F & F \\ \hline
            T & F & F & F & F & F & F \\ \hline
            F & T & T & F & T & F & F \\ \hline
            F & T & F & F & F & F & F \\ \hline
            F & F & T & F & F & F & F \\ \hline
            F & F & F & F & F & F & F \\ \hline
        \end{tabular}
    \end{table}
    From the above truth table, the result is established. We now show $p \lor (q \lor r) \equiv (p \lor q) \lor r$.
    \begin{table}[H]
        \centering
        \begin{tabular}{|c|c|c||c|c||c|c|}
            \hline
            $\boldsymbol{p}$ & $\boldsymbol{q}$ & $\boldsymbol{r}$ & $\boldsymbol{p \lor q}$ & $\boldsymbol{q \lor r}$ & $\boldsymbol{p \lor (q \lor r)}$ & $\boldsymbol{(p \lor q) \lor r}$ \\ \hhline{|=|=|=#=|=#=|=|}
            T & T & T & T & T & T & T \\ \hline
            T & T & F & T & T & T & T \\ \hline
            T & F & T & T & T & T & T \\ \hline
            T & F & F & T & F & T & T \\ \hline
            F & T & T & T & T & T & T \\ \hline
            F & T & F & T & T & T & T \\ \hline
            F & F & T & F & T & T & T \\ \hline
            F & F & F & F & F & F & F \\ \hline
        \end{tabular}
    \end{table}
    This proves the second claim.
\end{proof}

\begin{proposition}[Distributivity of AND and OR]
    Let $p$, $q$, and $r$ be statements. Then $p \land (q \lor r) \equiv (p \land q) \lor (p \land r)$ and $p \lor (q \land r) \equiv (p \lor q) \land (p \lor r)$.
\end{proposition}
\begin{proof}
    We first prove the first claim.
    \begin{table}[H]
        \centering
        \begin{tabular}{|c|c|c||c|c|c||c|c|}
            \hline
            $\boldsymbol{p}$ & $\boldsymbol{q}$ & $\boldsymbol{r}$ & $\boldsymbol{p \land q}$ & $\boldsymbol{p \land r}$ & $\boldsymbol{q \lor r}$ & $\boldsymbol{p \land (q \lor r)}$ & $\boldsymbol{(p \land q) \lor (p \land r)}$ \\ \hhline{|=|=|=#=|=|=#=|=|}
            T & T & T & T & T & T & T & T \\ \hline
            T & T & F & T & F & T & F & F \\ \hline
            T & F & T & F & T & T & T & T \\ \hline
            T & F & F & F & F & F & F & F \\ \hline
            F & T & T & F & F & T & F & F \\ \hline
            F & T & F & F & F & T & F & F \\ \hline
            F & F & T & F & F & T & F & F \\ \hline
            F & F & F & F & F & F & F & F \\ \hline
        \end{tabular}
    \end{table}
    We now prove the second claim.
    \begin{table}[H]
        \centering
        \begin{tabular}{|c|c|c||c|c|c||c|c|}
            \hline
            $\boldsymbol{p}$ & $\boldsymbol{q}$ & $\boldsymbol{r}$ & $\boldsymbol{p \lor q}$ & $\boldsymbol{p \lor r}$ & $\boldsymbol{q \land r}$ & $\boldsymbol{p \lor (q \land r)}$ & $\boldsymbol{(p \lor q) \land (p \lor r)}$ \\ \hhline{|=|=|=#=|=|=#=|=|}
            T & T & T & T & T & T & T & T \\ \hline
            T & T & F & T & T & F & T & T \\ \hline
            T & F & T & T & T & F & T & T \\ \hline
            T & F & F & T & T & F & T & T \\ \hline
            F & T & T & T & T & T & T & T \\ \hline
            F & T & F & T & F & F & F & F \\ \hline
            F & F & T & F & T & F & F & F \\ \hline
            F & F & F & F & F & F & F & F \\ \hline
        \end{tabular}
    \end{table}
    This proves the second claim.
\end{proof}

We now note an important set of properties regarding logical connectives.

\begin{proposition}[De Morgan's Laws]\label{prop:de-morgan}\index{De Morgan's Laws}
    Let $p$ and $q$ be statements. Then $\lnot (p \land q) \equiv (\lnot p) \lor (\lnot q)$ and $\lnot (p \lor q) \equiv (\lnot p) \land (\lnot q)$.
\end{proposition}
\begin{proof}
    See \myref{exercise:de-morgan} (later).
\end{proof}

We end this section by looking at the order of precedence for logical connectives.
\begin{table}[H]
    \centering
    \begin{tabular}{|c|c|}
        \hline
        \textbf{Connective} & \textbf{Precedence} \\ \hline
        $\lnot$ & 1 \\ \hline
        $\land$ and $\lor$ & 2 \\ \hline
        $\implies$ and $\iff$ & 3 \\ \hline
    \end{tabular}
\end{table}
To disambiguate the order of operations, we use parentheses (i.e., brackets).

\begin{exercise}\label{exercise:de-morgan}
    Prove \myref{prop:de-morgan}.
\end{exercise}

\begin{exercise}
    Let $p$, $q$, and $r$ be statements. Simplify the statement
    \[
        ((p \lor \lnot q) \land \lnot r) \lor p
    \]
    into a statement that uses only three logical connectives in total.
\end{exercise}

\section{Predicates and Quantifiers}
It would certainly be useful to include variables in our statements. Suppose I wanted to claim that ``if $n$ is odd, then $n + 1$ is not odd''. This looks like a conditional statement, but because $n$ is a variable, this is not a statement. In fact, such sentences are called predicates.

\begin{definition}[Predicates]
    A \term{predicate}\index{predicate} is a sentence that contains a finite number of \term{free variables}\index{free variable}\index{variable!free} and becomes a statement when specific values are substituted for the variables.
\end{definition}

For notational clarity, and to avoid confusing predicates with statements, we denote predicates using uppercase letters.

\begin{example}
    \begin{itemize}[leftmargin=*]
        \item Consider the sentence ``$n$ is a multiple of 3''. This is not a statement since we have the free variable ``$n$'' within the sentence; this is a predicate. Denoting this predicate by $P(n)$, we see that $P(0)$ is a true statement, $P(1)$ is a false statement, $P(2)$ is a false statement, $P(3)$ is a true statement, and so on.
        \item Let the predicate $Q(n)$ be ``$n$ is a factor of 9''. Then $Q(1)$, $Q(3)$, and $Q(9)$ are true statements.
    \end{itemize}
\end{example}

Returning back to our opening example, we could let $P(n)$ denote the predicate ``if $n$ is odd, then $n + 1$ is not odd''. One observes that, no matter what integer we substitute for $n$, the statement $P(n)$ is always true. What we really want to do with ``if $n$ is odd, then $n + 1$ is not odd'' is to quantify that this is true for all integers $n$.

There are two main quantifiers that we will look at here. The first is called the universal quantifier.

\begin{definition}[Universal Quantification]
    The \term{universal quantifier}\index{quantifier!universal} is $\forall$ and is read ``for all'' or ``every''.

    A statement $q$ of the form ``$\forall x, P(x)$'' where $P(x)$ is a predicate is called a \term{statement!universal}. The statement $q$ is true if and only if $P(x)$ is true for all values of $x$.
\end{definition}

\begin{example}
    \begin{itemize}[leftmargin=*]
        \item We can express the (false) statement that ``every number is greater than 0'' as ``$\forall x, x>0$''. Here the predicate $P(x)$ is ``$x>0$''. This is clearly false since 0 is not greater than 0.
        \item We can express the (true) statement that ``for every integer $n$, the integer $2n$ is even'' as ``$\forall x, ((n \text{ is an integer}) \implies (2n \text{ is even}))$''.
    \end{itemize}
\end{example}
\begin{remark}[For Those Who Know Sets]
    We can write ``for every integer $n$, the integer $2n$ is even'' as ``$\forall x \in \Z, (2n \text{ is even})$'' instead.
\end{remark}

The second quantifier is called the existential quantifier.

\begin{definition}[Existential Quantification]
    The \term{existential quantifier}\index{quantifier!existential} is $\exists$ and is read ``there exists'' or ``there is''.

    A statement $q$ of the form ``$\exists x, P(x)$'' where $P(x)$ is a predicate is called a \term{existential statement}\index{statement!existential}. The statement $q$ is true if and only if $P(x)$ is true for at least one value of $x$.
\end{definition}

\begin{example}
    We can express the (true) statement that ``there is a number greater than 0'' as ``$\exists x, x > 0$''. Here the predicate $P(x)$ is ``$x > 0$''. This is clearly true since 1 is greater than 0.
\end{example}

\begin{remark}[In Aid of Clarity]
    We can write ``$\exists x, P(x)$'' as ``$\exists x \text{ such that } P(x)$''. We can write ``such that'' as ``s.t.'' to make it more concise, i.e. ``$\exists x, P(x)$'' is the same as ``$\exists x \text{ s.t. } P(x)$''.
\end{remark}

It is critical to note that the truth or falsity of such quantified statements depends on the \term{domain of discourse}\index{domain!of discourse}. Consider the quantified statement ``$\forall x, \exists y \text{ s.t. } (y<x)$''. Certainly if we are considering all the real numbers this is a true statement! We could easily set $y = x - 1$ and we would be done here. However, over the positive integers, such a trick no longer works. If $x = 1$, then $y = 1 - 1 = 0$ is not a positive integer. So it is essential to clarify the domain of discourse when writing down quantified statements.

\begin{exercise}
    Suppose our domain of discourse is the positive integers. Convert the statement ``for all $n > 2$ the equation $x^n + y^n = z^n$ has no solution'' into symbolic notation using quantifiers and connectives.
\end{exercise}

% TODO: Add negation of quantified statements

\section{Hierarchy of Mathematical Results}
Mathematical statements often have a certain `tier' attached to them. We look at the hierarchy of some of these `tiers'.
\begin{itemize}
    \item \term{Definition}\index{definition}: A definition is a precise, formal, and abstract statement that introduce and clarify the meaning of mathematical terms, concepts, and objects.
    \item \term{Axiom}\index{axiom}: A statement that is assumed to be true to act as a starting point for further reasoning.
    \item \term{Proposition}\index{proposition}: A proposition can be thought of as a general proven result. ``\term{Claims}'' and ``\term{Observations}'' are equivalent to propositions.
    \item \term{Lemma}\index{lemma}: A lemma can be thought of as a `small' proven result that can help build other mathematical results. For example, Euclid's lemma that ``if a prime $p$ divides the product $ab$ of two integers $a$ and $b$, then $p$ must divide at least one of those integers $a$ or $b$'' is used in the proof of the Fundamental Theorem of Arithmetic.
    \item \term{Theorem}\index{theorem}: A theorem can be thought of as a `big' proven result. For example, the Fundamental Theorem of Arithmetic is core to arithmetic as it describes how integers can be uniquely decomposed into its prime factors.
    \item \term{Corollary}\index{corollary}: A corollary is said to be a `follow-up result' from a theorem. These results usually follow `very quickly' from a theorem or a proposition. For example, the AM-GM inequality is a corollary of the Pythagorean theorem.
\end{itemize}
